\begin{corollary}{4.7.3}\label{IV.4.7.3}
Let \(X\in\Ob\mathcal C\), and \(R\) an (M)-effective equivalence
relation in \(X\); suppose moreover that \(R\to X\times X\)
belongs to (N), where (N) satisfies (a) and
\((\mathrm f_{\mathcal T})\).
Then, for every \(Y\) as in 4.7.2, \(R_Y\to Y\times Y\)
also belongs to (N), and hence, by 4.7.1, one has:
\[
(Y/R_Y)(S)=\left\{\begin{array}{l}
\text{sub-}S\text{-objects }Z\text{ of }Y_S\text{ stable under }R_Y\times S,\\
\text{such that }Z\to Y_S\text{ belongs to (N) (then }Z\to X_S
\text{ does too),}\\
Z\to S\text{ is covering, and }R_Z\simeq Z\times_S Z
\end{array}\right\}.
\]
\end{corollary}

\sgasection{5}{Passage to the quotient and algebraic structures}
\sgasection{5.1}{Principal homogeneous bundles}
\begin{definition}{5.1.0}\label{IV.5.1.0}
\nde{The numbering 5.1.0 has been introduced for later references.}
Recall (III 0.1) that an object \(X\) with a group of operators
\(H\) (on the right) is said to be \emph{formally principal
homogeneous} under \(H\)\nde{One also says \og pseudo-torsor\fg,
cf. EGA IV$_4$, 16.5.15. Moreover, the more general notion
of formally homogeneous object (not necessarily principal
homogeneous) is defined in the addition 6.7.1 at the end of
this exposé.} if the canonical morphism (of functors)
\[
X\times H\longrightarrow X\times X
\]
defined by \((x,h)\mapsto(x,xh)\) is an isomorphism.
It amounts to the same to say (cf. \loccit) that for every
\(S\in\Ob\mathcal C\), \(X(S)\) is formally principal
homogeneous under \(H(S)\), that is, empty or principal
homogeneous under \(H(S)\).
In particular, if one lets \(H\) act on itself by translations
(on the right), \(H\) becomes formally principal homogeneous
under itself.
\end{definition}
\begin{definition}{5.1.1}\label{IV.5.1.1}
The object \(X\) with a group of operators \(H\) is said
to be \emph{trivial} if it is isomorphic (as an object with
a group of operators \(H\)) to \(H\) on which \(H\) acts
by translations.
\end{definition}
\begin{proposition}{5.1.2}\label{IV.5.1.2}
Let \(X\) be formally principal homogeneous under \(H\).
One has an isomorphism
\[
\Gamma(X)\isomto\Isom_{H\text{-obj.}}(H,X)
\]
of principal homogeneous sets under \(\Gamma(H)\).
\end{proposition}
\begin{proof}
To every section \(x\) of \(X\) one associates the morphism
from \(H\) to \(X\) defined set-theoretically by \(h\mapsto xh\).
The stated assertion is immediate, by reduction to the
set-theoretic case.
\end{proof}
\begin{corollary}{5.1.3}\label{IV.5.1.3}
One has an isomorphism of objects with operators \(H\)
\[
X\isomto\underline{\Isom}_{H\text{-obj.}}(H,X).
\]
\end{corollary}
\begin{corollary}{5.1.4}\label{IV.5.1.4}
For an object with a group of operators to be trivial,
it is necessary and sufficient that it be formally principal
homogeneous and that it possess a section.
\end{corollary}
\begin{definition}{5.1.5}\label{IV.5.1.5}
Let \(\mathcal C\) be a category endowed with a topology.
One says that the \(S\)-object \(X\) with an \(S\)-group
of operators \(H\) is a \emph{principal homogeneous bundle}
under \(H\) if it is locally trivial, that is, if the
following equivalent conditions hold:
\nde{In this case one also says that \(X\) is an \(H\)-torsor.}
\begin{enumeratei}
\item The set of \(T\to S\) such that (the functor)
\(X\times_S T\) is trivial under \(H\times_S T\) is
a refinement of \(S\).
\item There exists a covering family (for a pretopology
defining the given topology) \(\{S_i\to S\}\) such that
for each \(i\), the \(S_i\)-functor \(X\times_S S_i\)
with an \(S_i\)-group-functor of operators \(H\times_S S_i\)
is trivial (= possesses a section over \(S_i\)).
\end{enumeratei}
\end{definition}
\begin{proposition}{5.1.6}\label{IV.5.1.6}
Let \(\mathcal C\) be a category endowed with a topology
\(\mathcal T\). Let (M) be a family of morphisms of
\(\mathcal C\) satisfying axioms (a) to
\((\mathrm e_{\mathcal T})\) of 4.6.3.
Let \(H\) be an \(S\)-group such that the structural
morphism \(H\to S\) is an element of (M), and \(P\)
an \(S\)-object with an \(S\)-group of operators \(H\).
The following conditions are equivalent:
\begin{enumeratei}
\item \(P\) is a principal homogeneous bundle under
\(H\) (Definition 5.1.5).
\item \(P\) is formally principal homogeneous under
\(H\) and the structural morphism \(P\to S\) is
an element of (M).
\item There exists a morphism \(S'\to S\) which is an
element of (M), such that by base extension from \(S\)
to \(S'\), \(P\) becomes trivial, that is,
\(P\times_S S'\) is trivial under \(H\times_S S'\).
\item \(H\) acts freely on \(P\), in an (M)-effective
manner, and the quotient \(P/H\) is isomorphic to \(S\).
\end{enumeratei}
\end{proposition}
\begin{proof}
First remark that (ii) and (iv) are equivalent, taking
account of the fact that in either case \(P\to S\)
is an element of (M), hence base-changeable, which
ensures the representability of the fibered products
\(H\times_S P\), \(P\times_S P\).
It is clear that (ii) implies (iii), for one may take
\(P\) itself as \(S'\), the hypothesis that \(P\)
is formally principal homogeneous implying that
\(P\times_S P\) is trivial under \(H\times_S P\)
(5.1.4), for it has a section (the diagonal section).
It is clear that (iii) implies (i), for \(\{S'\to S\}\)
is a covering family by axiom \((\mathrm d_{\mathcal T})\).
It therefore remains to show that (i) implies (ii).
The morphism of sheaves \(P\times_S H\to P\times_S P\)
is locally an isomorphism, hence an isomorphism (4.5.8);
thus \(P\) is formally principal homogeneous.
The structural morphism \(P\to S\) is locally isomorphic
to the structural morphism \(H\to S\), which is an element
of (M). It is therefore itself an element of (M)
by \((\mathrm e_{\mathcal T})\).
\end{proof}
The equivalence between (i) and (iv) generalizes:
\begin{proposition}{5.1.7}\label{IV.5.1.7}
Under the same hypotheses on \(\mathcal C\) and (M),
let \(H\) be an \(S\)-group, and \(X\) an \(S\)-object
on which \(H\) acts (on the right).
Suppose that the structural morphism \(H\to S\)
is an element of (M). The following conditions are equivalent:
\begin{enumeratei}
\item \(H\) acts freely on \(X\), in an (M)-effective manner.
\item There exists an \(S\)-morphism \(p:X\to Y\)
compatible with the equivalence relation defined in
\(X\) by the action of \(H\), such that the action
of \(H\times_S Y\) on \(X\) over \(Y\) deduced from
it makes \(X\) a principal homogeneous bundle under
\(H_Y\) over \(Y\).
\end{enumeratei}
Under these conditions \(p\) identifies \(Y\) with
the quotient \(X/H\).
\end{proposition}
\begin{proof}
If \(p:X\to Y\) is a morphism compatible with the
action of \(H\), then the action of \(H\times_S Y\)
on \(X\) over \(Y\) deduced from it defines in \(X\)
the same equivalence relation as the action of \(H\),
by virtue of the formula
\[
H_Y\times_Y X\isomto H\times_S X.
\]
The proposition follows from this remark and the
preceding equivalence (iv) \(\Leftrightarrow\) (i).
\end{proof}
\begin{corollary}{5.1.7.1}\label{IV.5.1.7.1}
\nde{This particular case has been added as a corollary,
and will be used several times in the following exposés.}
Let \(\mathcal C\) be a category possessing a final object,
stable under fibered products and endowed with a
topology \(\mathcal T\) less fine than the canonical topology.
Let \(f:G\to H\) be a morphism of \(\mathcal C\)-groups,
and \(K=\Ker(f)\). Suppose that \(f\) is covering for
the topology \(\mathcal T\).

Then \(H\) represents the quotient sheaf \(G/K\),
and \(f\) is a \(K_H\)-torsor.
(N.B. One will also say: \og \(G\) is a \(K\)-torsor
over \(H\)\fg.)
\end{corollary}
\begin{proof}
Indeed, since \(f\) is covering, it is a universal
effective epimorphism (4.4.3), hence by 3.3.3.1,
\(H\) is the quotient of \(G\) by the equivalence
relation \(\mathcal R(f)=G\times_H G\).
On the other hand, the morphism \(G\times K\to G\times_H G\),
\((g,k)\mapsto(g,gk)\), is an isomorphism of objects
with a group of operators \(K_G=G\times_H K_H\)
(its inverse being given by \((g,g')\mapsto(g,g^{-1}g')\)).
Thus, on the one hand, \(\mathcal R(f)\) is the
equivalence relation defined by \(K\); on the other,
since the morphism \(f:G\to H\) is covering,
\(f\) is a \(K_H\)-torsor by 5.1.6 (ii)
(or directly by Definition 5.1.5 (ii)).
\end{proof}
We may now specify Theorem 4.6.6 in the case of passage
to the quotient by a group of operators:
\begin{proposition}{5.1.8}\label{IV.5.1.8}
Under the hypotheses of 5.1.7, denote by \(F_0\)
the functor over \(S\) defined as follows:
for each \(S'\to S\), \(F_0(S')\) is the set of
representable sub-\(S'\)-functors \(Z\) of \(X\times_S S'\),
stable under \(H\times_S S'\) and principal homogeneous
bundles under this \(S'\)-group for the induced action (3.2.2).

(i) The following conditions are equivalent:

a) The action of \(H\) on \(X\) is (M)-effective and free.
\nde{\og and free\fg{} has been added.}

b) \(F_0\) is representable.

Under these conditions \(F_0=X/H\).

(ii) Let (N) be a family of morphisms stable under
base change, such that for every covering family
\(\{S'_i\to S'\}\) and every family \(\{T_i\to S'_i\}\)
of morphisms of (N), every descent datum on the
\(T_i\) relative to \(\{S'_i\to S'\}\) is effective.
Suppose that the morphism \(X\times_S H\to X\times_S X\)
is an element of (N), and that \(X\) is base-changeable.
Then the element \(p\) of \(\Hom(X,F_0)\) corresponding
to the subobject \(X\times_S H\) of \(X\times_S X\)
identifies \(F_0\) with the quotient sheaf \(X/H\).
\end{proposition}

\sgasection{5.2}{Group structures and passage to the quotient}
In this number we concern ourselves with the algebraic
structures that one can put on the quotient \(G/H\)
of a group by a subgroup.
We shall first work in the category of sheaves on
\(\mathcal C\) for an arbitrary topology.
Taking the canonical topology and using 4.5.12,
we shall obtain results for passage to the universal
effective quotient in \(\mathcal C\).

\begin{proposition}{5.2.1}\label{IV.5.2.1}
Let \(u:H\to G\) be a monomorphism of sheaves of groups.
There exists on the quotient sheaf \(G/H\) a unique
structure of object with a group of operators \(G\)
such that the canonical morphism
\[
p:G\longrightarrow G/H
\]
is a morphism of objects with a group of operators \(G\).
This structure is functorial in the pair \((G,H)\):
if one has a commutative diagram
\[
\xymatrix{H\ar[r]\ar[d] & G\ar[d]^f\\ H'\ar[r] & G' },
\]
the morphism \(\overline f:G/H\to G'/H'\) (3.2.3)
is compatible with the morphism \(f\) on the groups
of operators.
\end{proposition}
\begin{proof}
Indeed, the sheaf \(G/H\) is the sheaf associated to the presheaf
\[
i(G)/i(H):S\longmapsto G(S)/H(S);
\]
since the functor \(a\) is left exact, it transforms
objects with groups of operators into objects with
a group of operators.
Since the presheaf \(i(G)/i(H)\) is endowed with a
structure of object with a group of operators \(i(G)\),
\(G/H=a(i(G)/i(H))\) is endowed with a structure of
object with operators \(a(i(G))=G\).
This structure evidently enjoys all the stated properties.
\end{proof}
\begin{corollary}{5.2.2}\label{IV.5.2.2}
Let \(u:H\to G\) be a monomorphism of \(\mathcal C\)-groups.
Suppose that the action of \(H\) on \(G\) is universal
effective. There exists on the quotient object
\(G/H\in\Ob\mathcal C\) a unique structure of object
with a group of operators \(G\) such that \(p:G\to G/H\)
is a morphism of objects with operators.
This structure is functorial in the pair \((H,G)\)
(\(H\) acting in a universal effective manner on \(G\))
in the preceding sense.
\end{corollary}
\begin{proposition}{5.2.3}\label{IV.5.2.3}
Let \(u:H\to G\) be a monomorphism of sheaves of groups
identifying \(H\) with an invariant subsheaf of groups
of \(G\). There exists on the quotient sheaf \(G/H\)
a unique structure of sheaf of groups such that
the canonical morphism \(p:G\to G/H\) is a morphism
of groups. This structure is functorial in the pair
\((H,G)\) (\(H\) invariant).
\end{proposition}
\begin{proof}
The proof is similar to that of 5.2.1.
\end{proof}
\begin{corollary}{5.2.4}\label{IV.5.2.4}
Let \(u:H\to G\) be a monomorphism of \(\mathcal C\)-groups
identifying \(H\) with an invariant subgroup of \(G\).
Suppose that the action of \(H\) on \(G\) is universal
effective. There exists on the quotient object
\(G/H\in\Ob\mathcal C\) a unique group structure such
that the canonical morphism \(G\to G/H\) is a morphism
of groups. This structure is functorial in the pair
\((H,G)\) (\(H\) invariant, \(H\) acting in a universal
effective manner).
\end{corollary}
One can characterize the group structure of \(G/H\)
in a more expressive manner:
\begin{proposition}{5.2.5}\label{IV.5.2.5}
Under the conditions of 5.2.4, let \(K\) be a
\(\mathcal C\)-group, and \(f:G\to K\) a morphism.
The following conditions are equivalent:
\begin{enumeratei}
\item \(f\) is a morphism of groups compatible with
the equivalence relation defined by \(H\).
\item \(f\) is a morphism of groups inducing the
trivial morphism \(H\to K\).
\item \(f\) factors through a morphism of groups \(G/H\to K\).
\end{enumeratei}
In particular, one has an isomorphism functorial in the group \(K\)
\[
\Hom_{\mathcal C\text{-gr.}}(G/H,K)
\isomto\{f\in\Hom_{\mathcal C\text{-gr.}}(G,K)\mid f\circ u=e\}.
\]
\end{proposition}
\begin{proof}
The equivalence of (i) and (ii) is proved set-theoretically.
One evidently has (iii) \(\Rightarrow\) (ii).
The equivalence of (iii) and (ii) follows from the formula
\[
\Hom(G/H,K)\simeq\Hom(i(G)/i(H),K)
\]
and the definition of the group structure of \(G/H\).
\end{proof}
\begin{remark}{5.2.6}\label{IV.5.2.6}
In the preceding situation, if the kernel of \(f\)
is exactly \(H\), the morphism \(G/H\to K\) through
which \(f\) factors is a monomorphism.
This follows at once from 3.3.4.
\end{remark}
In the case of sheaves of groups, one can specify 4.4.11 by the
\begin{proposition}{5.2.7}\label{IV.5.2.7}
Let \(G\) be a sheaf of groups and \(H\) an invariant
subsheaf of groups. For every subsheaf of groups \(K\)
of \(G\) containing \(H\), let \(K'\) be the quotient
group \(K/H\) considered as a subgroup of \(G'=G/H\).

One has \(K=K'\times_{G'}G\), and the maps
\(K\mapsto K/H\), \(K'\mapsto K'\times_{G'}G\)
establish a bijection between the set of subsheaves
of groups of \(G\) containing \(H\) and the set of
subsheaves of groups of \(G'\).
Under this correspondence, invariant subsheaves of
groups of \(G,G'\) correspond.
\end{proposition}
\begin{proof}
The first part follows easily from 4.4.11 and 3.2.4.
It remains to see that \(K\) is invariant in \(G\)
if and only if \(K'\) is invariant in \(G'\).
If \(K\) is invariant in \(G\), then the presheaf
\(i(K)/i(H)\) is invariant in \(i(G)/i(H)\).
The same is true of the associated sheaves, by
virtue of the customary argument.
Conversely, if \(K'\) is invariant in \(G'\), then
the fibered product \(K\times_{G'}G\) is invariant
in \(G\), as one sees immediately.
% typo?: kept as printed: K rather than K' in the last fibered product.
\end{proof}
If now \(L\) is any subsheaf of groups of \(G\),
let \(\overline L\) be the saturation of \(L\) for
the equivalence relation defined by \(H\);
one will also write \(\overline L=L\cdot H\).
\begin{proposition}{5.2.8}\label{IV.5.2.8}
Under the preceding conditions, \(L\cdot H\) is a
subsheaf of groups of \(G\) containing \(H\), and
the image of \(L\) in \(G/H\) is identified with
\[
(L\cdot H)/H\simeq L/(H\cap L).
\]
\end{proposition}
\begin{proof}
Indeed, denote by \(L'\) the image sheaf of \(L\)
in \(G/H\). It is a subsheaf of groups of \(G/H\)
corresponding to \(L\cdot H\) under the correspondence
of the preceding proposition.
Since the morphism \(L\to L'\) is covering, hence
a universal effective epimorphism of sheaves,
it follows from 4.4.9 that \(L'\) is identified
with the quotient of \(L\) by the kernel of
\(L\to L'\), which is evidently precisely \(H\cap L\).
\end{proof}

Finally consider the following situation:
one has a sheaf of groups \(G\), a subsheaf of
groups \(K\), and a subsheaf of groups \(H\) of
\(K\), invariant in \(K\).
First let us define an action (on the right) of
the sheaf of groups \(H\backslash K\) (\(=K/H\))
on \(G/H\). The group \(K\) acts by right translations
on \(G\). Since \(H\) is invariant in \(K\), this
action is compatible with the equivalence relation
defined by the action of \(H\), and therefore defines
an action of \(K\) on \(G/H\), that is, a morphism
from the group \(K^\circ\) opposite to \(K\)
to \(\underline{\Aut}(G/H)\).
Since the latter is a sheaf (4.5.13) and this
morphism is trivial on \(H\), it factors through
\(K/H\) and defines the desired action.
Since the right and left actions of \(G\) on
itself commute, the actions of \(G,K/H\) on
\(G/H\) commute.

\begin{proposition}{5.2.9}\label{IV.5.2.9}
Under the preceding conditions, \(K/H\) acts
freely (on the right) on \(G/H\), and one has
a canonical isomorphism of sheaves with a
group of operators \(G\)
\[
(G/H)/(K/H)\simeq G/K.
\]
When \(K\) is invariant in \(G\), in which case
\(K/H\) is invariant in \(G/H\) (5.2.7), this
isomorphism respects the group structures of
the two members.
\end{proposition}
\begin{proof}
One has an isomorphism of presheaves
\[
(i(G)/i(H))/(i(K)/i(H))\isomfrom i(G)/i(K),
\]
which respects the structures of objects with
a group of operators \(i(G)\).
The announced result is obtained by applying
the functor \(a\) to this relation.
\end{proof}
\begin{corollary}{5.2.10}\label{IV.5.2.10}
Let \(G\) be a \(\mathcal C\)-group, \(K\) a
sub-\(\mathcal C\)-group of \(G\), and \(H\)
an invariant sub-\(\mathcal C\)-group of \(K\).
Let (M) be a family of morphisms of \(\mathcal C\)
satisfying axioms (a) to \((\mathrm e_{\mathcal T})\).
Suppose that the action of \(H\) on \(G\)
(resp. \(K\)) on the right is (M)-effective.
Then \(K/H\) acts naturally and freely on
the right on \(G/H\); this action commutes
with that of \(G\). The following conditions are equivalent:
\begin{enumeratei}
\item The action of \(K\) on \(G\) is (M)-effective.
\item The action of \(K/H\) on \(G/H\) is (M)-effective.
\end{enumeratei}
Under these conditions one has an isomorphism
of objects\nde{Of \(\mathcal C\).} with a group of operators \(G\):
\[
(G/H)/(K/H)\simeq G/K.
\]
\end{corollary}
